\chapter{Vigtige begreber}
\label{chap:ic}

Dette bilag er en samling af vigtige matematiske begreber, som ofte bruges i softwareudvikling. 
Indholdet i bilaget bygger på begreberne i denne bog.

% ---------- 01 

\begin{proposition}[Regneregler for brøker]
For $a,b,c,m \in \mathbb{R}$, hvor $a,b,c,m \neq 0$ når det er nødvendigt, gælder følgende identiteter:
\begingroup
\setlength{\jot}{8pt} % increase space between align rows (default ~3pt)
\begin{align*}
(1) & \quad \frac{a}{b} \times m = \frac{am}{b} \\
(2) & \quad \frac{a}{b} \div m = \frac{a}{bm} \\
(3) & \quad m \div \frac{a}{b} = \frac{mb}{a} \\
(4) & \quad \frac{a}{b} \times \frac{c}{a} = \frac{c}{b} \\
(5) & \quad \frac{a}{b} \div \frac{c}{a} = \frac{a^2}{bc} \\
(6) & \quad \frac{a}{b} = \frac{ac}{bc} \\
(7) & \quad \frac{a}{b} + \frac{c}{a} = \frac{a^2 + bc}{ab}
\end{align*}
\endgroup
\end{proposition}

\begin{proposition}{Regneregler for rødder}
    \begin{align*}
    &(1) \quad \sqrt[n]{x y}=\sqrt[n]{x} \times \sqrt[n]{y}   &   &\textnormal{begge sider reelle; for lige }n\textnormal{ kræves }x,y\geq0\\
    &(2) \quad \sqrt[n]{\frac{x}{y}} = \frac{\sqrt[n]{x}}{\sqrt[n]{y}}    &   &\textnormal{begge sider reelle, }y\neq0;\textnormal{ for lige }n\textnormal{ kræves }x\geq0,y>0 \\
    &(3) \quad \sqrt{x^2} = |x|   &   &\textnormal{hvor } x \in \mathbb{R} \\
    &(4) \quad (\sqrt[n]{x})^n = x    &   &\textnormal{for ulige }n\textnormal{ og alle }x;\ \textnormal{for lige }n\textnormal{ og }x\geq0\\
    &(5) \quad \sqrt[n]{-x}=-\sqrt[n]{x}  &   &\textnormal{hvor } x \geq 0 \textnormal{ og } n \in \mathbb{N} \textnormal{ er ulige }
    \end{align*}
\end{proposition}

\begin{proposition}{Potensregneregler for heltallige eksponenter}
Lad $n,m\in\mathbb{Z}$. Følgende identiteter gælder, når potenserne er definerede. Især skal et grundtal være forskelligt fra nul, når der optræder en negativ eksponent:
\[
\begin{aligned}
&\text{(1)}\quad x^{n}\cdot x^{m}=x^{\,n+m},\\[2pt]
&\text{(2)}\quad \dfrac{x^{n}}{x^{m}}=x^{\,n-m}\quad \text{med } x\neq 0,\\[2pt]
&\text{(3)}\quad x^{n}\cdot y^{n}=(xy)^{n},\\[2pt]
&\text{(4)}\quad \dfrac{x^{n}}{y^{n}}=\left(\dfrac{x}{y}\right)^{n}\quad \text{med } y\neq 0,\\[2pt]
&\text{(5)}\quad \bigl(x^{n}\bigr)^{m}=x^{\,nm},\\[2pt]
&\text{(6)}\quad x^{1}=x.
\end{aligned}
\]
\end{proposition}

\begin{proposition}{Flere potensregneregler for heltallige eksponenter}
Lad $n,m\in\mathbb{Z}$. Så gælder følgende (med $x,y\in\mathbb{R}$ og forskellige fra nul, hvor det er angivet):
\[
\begin{aligned}
&\text{(7)}\quad x^{0}=1   &\qquad&   x \neq 0  \\
&\text{(8)}\quad \frac{1}{x^{m}}=x^{-m} &\qquad& x \neq 0 \\
\end{aligned}
\]
\end{proposition}

\begin{definition}{Det udvidede potensbegreb}
$\textnormal{Lad } m \in \mathbb{Z} \textnormal{ og lad }  n \in \mathbb{N}\textnormal{, så } \dfrac{m}{n} \in \mathbb{Q}  \textnormal{. Så gælder følgende:}$

\[
x^{\frac{m}{n}} = \sqrt[n]{x^m} \hspace{1cm} x > 0
\]

\textnormal{Mere specifikt gælder}

\[
x^{\frac{1}{n}} = \sqrt[n]{x} \hspace{1.3cm} x > 0
\]

\end{definition}

\begin{custombox}{Regler for omskrivning af formler}
Følgende operationer kan udføres på begge sider af formlen:
\begin{itemize}
    \item Læg den samme størrelse til på begge sider
    \item Træk den samme størrelse fra på begge sider
    \item Gang begge sider med den samme størrelse -- husk at gange alle led
    \item Divider begge sider med den samme størrelse forskellig fra nul -- husk at dividere alle led
    \item Anvend en funktion på begge sider. Det bevarer implikationen, men det bevarer kun ækvivalensen, når funktionen er injektiv på den relevante definitionsmængde. Kvadrering kan give falske løsninger, og at tage den reciprokke værdi kræver, at siderne er forskellige fra nul.
\end{itemize}
\end{custombox}

\begin{definition}
Lad $A$ og $B$ være ikke-tomme mængder. En funktion $f$ fra $A$ til $B$ er en tilordning, der til hvert element i $A$ knytter præcis ét element i $B$. Vi skriver $f(a)=b$, hvis $b$ er det entydige element i $B$, som funktionen $f$ knytter til elementet $a$ i $A$. Hvis $f$ er en funktion fra $A$ til $B$, skriver vi $f: A \rightarrow B$.
\end{definition}

\begin{definition}
Hvis $f$ er en funktion fra $A$ til $B$, siger vi, at $A$ er \textbf{definitionsmængden} for $f$, og at $B$ er \textbf{dispositionsmængden} for $f$. Hvis $f(a)=b$, siger vi, at $b$ er \textbf{billedet} af $a$, og at $a$ er et \textbf{urbillede} af $b$. \textbf{Værdimængden}, eller billedet, af $f$ er mængden af alle billeder af elementer i $A$. Hvis $f$ er en funktion fra $A$ til $B$, siger vi desuden, at $f$ \textbf{afbilder} $A$ ind i $B$.    
\end{definition}

\begin{definition} {Injektive funktioner}
En funktion \( f: A \rightarrow B \) kaldes \textbf{injektiv} (på engelsk \textbf{one-to-one}), hvis forskellige elementer i \( A \) afbildes over i forskellige elementer i \( B \). Med andre ord: Hvis \( f(a_1) = f(a_2) \), så er \( a_1 = a_2 \). Denne egenskab sikrer, at to forskellige elementer i \( A \) aldrig afbildes over i det samme element i \( B \).     
\end{definition}

\begin{definition}{Surjektive funktioner}
En funktion \( f: A \rightarrow B \) kaldes \textbf{surjektiv} (på engelsk \textbf{onto}), hvis ethvert element i \( B \) er billedet af mindst ét element i \( A \). Med andre ord findes der for ethvert \( b \in B \) mindst ét \( a \in A \), så \( f(a) = b \). Denne egenskab sikrer, at funktionen »dækker« hele mængden \( B \).
\end{definition}

\begin{definition}{Inverse funktioner}
Lad $f$ være en bijektion fra mængden $A$ til mængden $B$. Den inverse funktion til $f$ er den funktion, der til et element $b$ i $B$ knytter det entydige element $a$ i $A$, for hvilket $f(a)=b$. Den inverse funktion til $f$ betegnes $f^{-1}$. Altså er $f^{-1}(b)=a$, når $f(a)=b$.    
\end{definition}

\begin{definition}
    

Lad \( f: B \rightarrow C \) og \( g: A \rightarrow B \) være to funktioner. Den \textbf{sammensatte funktion} af \( f \) og \( g \), betegnet \( f \circ g \), er en funktion fra \( A \) til \( C \), der er defineret ved

\[
(f \circ g)(x) = f(g(x)),
\]

for ethvert \( x \in A \).

\end{definition}

\begin{definition}{Særlige eksponentialfunktioner}

\begin{align*}
y &= a \cdot 10^{bx} \hspace{1cm} \text{eksponentialfunktion med grundtal 10} \\
y &= a \cdot e^{bx}  \hspace{1.2cm} \text{naturlig eksponentialfunktion (grundtal }e\text{),}
\end{align*}

hvor $a$ og $b$ er konstanter, og definitionsmængden er alle reelle tal.
    
\end{definition}

\begin{definition}{Titalslogaritmer}

\[
\log x=y \iff 10^y=x, \qquad x>0
\]

\textit{Med ord}: $\log x$ er den eksponent i potensen af 10, der giver $x$   
\end{definition}

\begin{custombox}{Regneregler for titalslogaritmer}
Følgende identiteter kræver $x>0$ og $y>0$:
\begin{itemize}
    \item Logaritmen til et produkt:
    
    \[
    \log(xy)=\log x+\log y
    \]
    
    \textit{Med ord}: $\log$ til et produkt er lig med summen af logaritmerne til faktorerne.
    \vspace{0.2cm}
    \item Logaritmen til en kvotient:
    
    \[
    \log \frac{x}{y}=\log x-\log y
    \]
    \vspace{0.1cm}
    \textit{Med ord}: $\log$ til en kvotient er lig med logaritmen til tælleren minus $\log$ til nævneren.
    \vspace{0.2cm}
    \item Logaritmen til en potens:
    \[
    \log x^y=y \log x
    \]
    \textit{Med ord}: $\log$ til en potens er lig med eksponenten gange logaritmen til grundtallet.
\end{itemize}

   
\end{custombox}

\begin{custombox}{Det vigtigste at huske om logaritmer er dette}
    \textbf{En logaritme er en eksponent.}
\end{custombox}

\begin{definition}{Logaritme med vilkårligt grundtal}
\textit{Algebraisk}:

\hspace{1cm}$\log _b x=y$ hvis og kun hvis $b^y=x, \quad$ hvor $b>0, b \neq 1$, og $x>0$

\vspace{0.3cm}

\textit{Med ord}:

\hspace{1cm} $\log _b x=y$ betyder, at $y$ er den eksponent til $b$, der giver $x$ som resultat.
    
\end{definition}

\begin{definition}{Titalslogaritme og naturlig logaritme}
\hspace{1cm} \textit{Titalslogaritme}: Symbolet $\log x$ betyder $\log _{10} x$.

\hspace{1cm}  \textit{Naturlig}: \hspace{0.2cm}Symbolet $\ln x$ betyder $\log _e x$, hvor $e$ er en konstant lig med $2.71828182845 \ldots$
\end{definition}

\begin{custombox}{Reglen for basisskift for logaritmer}

\begin{equation*}
\log _a x=\frac{\log_b x}{\log_b a} \quad \text { eller } \quad \log_a x=\frac{1}{\log_b a}\left(\log_b x\right)
\end{equation*}
hvor $a,b,x>0$ og $a\neq1$, $b\neq1$.
    
\end{custombox}

\begin{custombox}{Logaritmeregneregler}
\setlength{\leftskip}{1cm}  % Reset the text indent
\setlength{\rightskip}{1cm} % Reset the right indent
Vi antager overalt, at logaritmernes grundtal er positive og forskellige fra 1, og at alle logaritmernes argumenter er positive.
Logaritmen til en potens:
$$
\log _b x^y=y \log _b x
$$

\textit{Med ord}: Logaritmen til en potens er lig med produktet af eksponenten og logaritmen til grundtallet.
\vspace{0.5cm}
Logaritmen til et produkt:
$$
\log _b(x y)=\log _b x+\log _b y
$$
\textit{Med ord}: Logaritmen til et produkt er lig med summen af logaritmerne til faktorerne.
\vspace{0.5cm}
Logaritmen til en kvotient:
$$
\log _b \frac{x}{y}=\log _b x-\log _b y
$$

\textit{Med ord}: Logaritmen til en kvotient er lig med logaritmen til tælleren minus logaritmen til nævneren.

\setlength{\leftskip}{0cm}  % Reset the text indent
\setlength{\rightskip}{0cm} % Reset the right indent

\end{custombox}

% ---------- 02

\begin{table}[ht]
\centering
\renewcommand{\arraystretch}{1.4}
\begin{tabular}{|c|c|c|c|}
\hline
\textbf{Talsystem} & \textbf{Symboler} & \textbf{Grundtal} & \textbf{Yderligere information} \\ \hline
\textbf{Decimalt} & 0-9 & 10 & - \\ \hline
\textbf{Binært} & 0, 1 & 2 & - \\ \hline
\textbf{Hexadecimalt} & 0-9, A-F & 16 & $\mathrm{A} \equiv 10, \mathrm{B} \equiv 11, \mathrm{C} \equiv 12,$ $\mathrm{D} \equiv 13, \mathrm{E} \equiv 14, \mathrm{F} \equiv 15$ \\ \hline
\textbf{Oktalt} & 0-7 & 8 & - \\ \hline
\end{tabular}
\caption{Oversigt over almindelige talsystemer}
\end{table}

\begin{table}[!ht]
\centering
\renewcommand{\arraystretch}{1.4}
\begin{tabular}{|c c c c c c c c|}
\hline 
\multirow{2}{*}{Decimalt tal} & & \multirow{2}{*}{Som potenser af 2} & \multicolumn{4}{c}{Potens af 2} & \multirow{2}{*}{Binært tal} \\
& & & 3 & 2 & 1 & 0 & \\ \hline
8 & $=$ & $2^3$ & 1 & 0 & 0 & 0 & 1000 \\
7 & $=$ & $2^2 + 2^1 + 2^0$ & 0 & 1 & 1 & 1 & 111 \\
6 & $=$ & $2^2 + 2^1$ & 0 & 1 & 1 & 0 & 110 \\
5 & $=$ & $2^2 + 2^0$ & 0 & 1 & 0 & 1 & 101 \\
4 & $=$ & $2^2$ & 0 & 1 & 0 & 0 & 100 \\
3 & $=$ & $2^1 + 2^0$ & 0 & 0 & 1 & 1 & 11 \\
2 & $=$ & $2^1$ & 0 & 0 & 1 & 0 & 10 \\
1 & $=$ & $2^0$ & 0 & 0 & 0 & 1 & 1 \\ \hline
\end{tabular}
\caption{Decimale tal i binær repræsentation}
\end{table}

\begin{proposition}{Regler for binær addition}
\[
0 + 0 = 0, \quad 0 + 1 = 1, \quad 1 + 0 = 1, \quad 1 + 1 = 10
\]
\end{proposition}

\begin{proposition}{Regler for binær multiplikation}
\[
\begin{aligned}
& 0 \times 0=0 \\
& 0 \times 1=0 \\
& 1 \times 0=0 \\
& 1 \times 1=1 \\
& 1 \times 10_2 = 10_2 \quad \left( \text{multiplikation med grundtallet } 10_2 \text{ tilføjer et 0 til sidst} \right)
\end{aligned}
\]
\end{proposition}

\begin{proposition}{XOR-operationen}
XOR giver 1, hvis de to bit, der sammenlignes, er forskellige, og 0, hvis de er ens:
\[
0 \oplus 0 = 0, \quad 0 \oplus 1 = 1, \quad 1 \oplus 0 = 1, \quad 1 \oplus 1 = 0
\]
\end{proposition}

\begin{figure}[htbp]
\centering
\begin{tikzpicture}[>=Stealth, node distance=4cm, thick, scale=1]

    % Draw larger circles for sets A and B
    \draw[thick] (0,0) circle [radius=1.5cm] node at (0, 2) {\Large $A$};
    \draw[thick] (6,0) circle [radius=1.5cm] node at (6, 2) {\Large $B$};

    % Place larger nodes for elements a and b=f(a)
    \node[fill, circle, inner sep=2pt, label=below:{$a$}] (a) at (0,0.3) {};
    \node[fill, circle, inner sep=2pt, label=below:{$b=f(a)$}] (b) at (6,0.3) {};

    % Draw the larger straight arrow representing f
    \draw[->, ultra thick, black] (a) -- (b) node[midway, above] {\Large $f$};

    % Draw the curved arrow representing f at the bottom, touching the circles
    \draw[->, ultra thick, black, bend right=30] ([shift={(270:1.5cm)}]0,0) 
    to node[midway, below] {\Large $f$} 
    ([shift={(270:1.5cm)}]6,0);

\end{tikzpicture}
\caption{En funktion \(f\), der afbilder et element \(a\) fra mængden \(A\) over i et element \(b=f(a)\) i mængden \(B\).}
\end{figure}

\begin{figure}[htbp]
\centering
\begin{tikzpicture}[>=Stealth, thick, scale = 1.5]

% Draw circles representing sets A and B
\draw (0,0) circle (1.5cm);
\draw (6,0) circle (1.5cm);

% Labels for sets A and B
\node at (0,-1) {A};
\node at (6,-1) {B};


% Points in sets with adjusted labels
\node[fill=black, circle, inner sep=1.5pt, label={[xshift=-0.6cm]below:{$a = f^{-1}(b)$}}] (A1) at (0,0.7) {};
\node[fill=black, circle, inner sep=1.5pt, label={[xshift=0.5cm]below:{$b = f(a)$}}] (B1) at (6,0.7) {};


% Arrows for f and f^{-1}
\draw[->] (A1) to[bend left=30] node[midway, above] {$f(a)$} (B1);
\draw[->] (B1) to[bend left=30] node[midway, below] {$f^{-1}(b)$} (A1);

% Arrows for f and f^{-1} below the main circles
\draw[->] (0.6,-1.5) to[bend left=20] node[midway, above] {$f$} (5.4,-1.5);
\draw[->] (5.4,-1.8) to[bend left=20] node[midway, below] {$f^{-1}$} (0.6,-1.8);

\end{tikzpicture}
\caption{Funktionen $f^{-1}$ er den inverse til funktionen $f$.}
\end{figure}

\begin{figure}[htbp]
\centering
\begin{tikzpicture}[>=Stealth, thick]

% Draw circles representing sets A, B, and C
\draw (0,0) circle (1.5cm);  % Set A
\draw (5,0) circle (1.5cm);  % Set B
\draw (10,0) circle (1.5cm); % Set C

% Labels for sets A, B, and C
\node at (0,-1) {A};
\node at (5,-1) {B};
\node at (10,-1) {C};

% Points in sets with adjusted labels
\node[fill=black, circle, inner sep=1.5pt, label={[xshift=-0.1cm]below:{$a$}}] (A) at (0,0.7) {};

\node[fill=black, circle, inner sep=1.5pt, label=below:{$g(a)$}] (B) at (5,0.7) {};
\node[fill=black, circle, inner sep=1.5pt, label={[xshift=0.4cm]below:{$f(g(a))$}}] (C) at (10,0.7) {};

% Arrows for g and f
\draw[->] (A) to[bend left=15] node[midway, above] {$g(a)$} (B);
\draw[->] (B) to[bend left=15] node[midway, above] {$f(g(a))$} (C);
\draw[->] ([shift={(360:1.5cm)}]0,0) to node[midway, below] {$g$} (3.5,0);
\draw[->] ([shift={(360:0cm)}]6.5,0) to node[midway, below] {$f$} (8.5,0);

% Arrows for f ∘ g and (f ∘ g)(a)
\draw[->] (A) to[bend left=40] node[midway, above] {$(f \circ g)(a)$} (C);
% Draw the curved arrow representing f at the bottom, touching the circles
\draw[->, bend right=20] ([shift={(270:1.5cm)}]0,0) to node[midway, below] {$f(g(a))$} ([shift={(270:0.7cm)}]10,0);

\end{tikzpicture}
\caption{Sammensætningen af funktionerne \(f\) og \(g\), betegnet \(f \circ g\), er den funktion, der fremkommer ved først at anvende \(g\) og derefter \(f\).}
\end{figure}


% ---------- 03
\begin{definition}[Delmængde]
    En mængde $A$ er en \textbf{delmængde} af en mængde $B$, hvis ethvert element i $A$ også er et element i $B$. Det skriver vi som $A \subseteq B$.
\end{definition}

\begin{definition}[Ægte delmængde]
    En mængde $A$ er en \textbf{ægte delmængde} af en mængde $B$, hvis $A \subseteq B$ og $A \neq B$. Det betyder, at $B$ skal indeholde mindst ét element, som ikke ligger i $A$. Det skriver vi som $A \subset B$.
\end{definition}

\begin{definition}[Universalmængde]
    \textbf{Universalmængden}, betegnet $U$, er mængden af alle mulige elementer, der betragtes i en given sammenhæng. Alle andre mængder i den sammenhæng betragtes som delmængder af universalmængden.
\end{definition}

\begin{definition}[Den tomme mængde]
    \textbf{Den tomme mængde} (eller \textbf{nulmængden}) er den entydige mængde, der ikke indeholder nogen elementer. Den betegnes $\emptyset$ eller $\{\}$.
\end{definition}

\begin{definition}[Disjunkte mængder]
    To mængder $A$ og $B$ er \textbf{disjunkte}, hvis de ikke har nogen elementer til fælles. Med andre ord er deres fællesmængde den tomme mængde: $A \cap B = \emptyset$.
\end{definition}

\begin{definition}[Kardinalitet]
    \textbf{Kardinaliteten} af en endelig mængde $A$, betegnet $|A|$, er antallet af forskellige elementer i mængden.
\end{definition}

\begin{definition}[Potensmængde]
    \textbf{Potensmængden} af en mængde $A$, betegnet $\mathcal{P}(A)$, er mængden af alle delmængder af $A$. Potensmængdens elementer er altså selv mængder.
\end{definition}

\begin{custombox}{Dualiteten mellem mængdelære og boolesk algebra}
    For en fast universalmængde er algebraen af dens delmængder med foreningsmængde, fællesmængde og komplement en boolesk algebra. Operationerne svarer til logisk OR, AND og NOT som følger:
    \begin{center}
    \renewcommand{\arraystretch}{1.3}
    \begin{tabular}{lcl}
        \textbf{Mængdelære} & & \textbf{Boolesk algebra / logik} \\
        \hline
        Foreningsmængde ($\cup$) & $\iff$ & Boolesk sum ($+$) eller \texttt{OR} \\
        Fællesmængde ($\cap$) & $\iff$ & Boolesk produkt ($\cdot$) eller \texttt{AND} \\
        Komplement ($A^c$) & $\iff$ & Komplementering ($\overline{x}$) eller \texttt{NOT} \\
        Universalmængden ($U$) & $\iff$ & Sand (1) \\
        Den tomme mængde ($\emptyset$) & $\iff$ & Falsk (0) \\
    \end{tabular}
    \end{center}
\end{custombox}

\begin{definition}[Fællesmængde]
    \textbf{Fællesmængden} af mængderne $A$ og $B$, betegnet $A \cap B$, er mængden af alle elementer, der ligger i både $A$ og $B$.
    \[ A \cap B = \{ x \mid x \in A \text{ og } x \in B \} \]
\end{definition}

\begin{definition}[Foreningsmængde]
    \textbf{Foreningsmængden} af mængderne $A$ og $B$, betegnet $A \cup B$, er mængden af alle elementer, der ligger i $A$, i $B$ eller i begge.
    \[ A \cup B = \{ x \mid x \in A \text{ eller } x \in B \} \]
\end{definition}

\begin{definition}[Mængdedifferens]
    \textbf{Differensen} mellem mængden $A$ og mængden $B$, betegnet $A \setminus B$, er mængden af alle elementer, der ligger i $A$, men ikke i $B$.
    \[ A \setminus B = \{ x \mid x \in A \text{ og } x \notin B \} \]
\end{definition}

\begin{definition}[Symmetrisk differens]
    Den \textbf{symmetriske differens} af mængderne $A$ og $B$, betegnet $A \oplus B$, er mængden af de elementer, der ligger i en af mængderne, men ikke i deres fællesmængde.
    \[ A \oplus B = (A \cup B) \setminus (A \cap B) \]
\end{definition}

\begin{definition}[Komplement]
    \textbf{Komplementet} til en mængde $A$, betegnet $A^c$, er mængden af alle elementer i universalmængden $U$, som ikke ligger i $A$.
    \[ A^c = U \setminus A \]
\end{definition}

\begin{definition}[Kartesisk produkt]
    Det \textbf{kartesiske produkt} af mængderne $A$ og $B$, betegnet $A \times B$, er mængden af alle mulige ordnede par $(a, b)$, hvor det første element $a$ er fra $A$, og det andet element $b$ er fra $B$.
    \[ A \times B = \{ (a, b) \mid a \in A \text{ og } b \in B \} \]
\end{definition}

\begin{custombox}{Mængdeoperationer som bitvise operationer}
    Lad bitstrengene for mængderne $A$ og $B$ være $s_A$ og $s_B$.
    \begin{itemize}
        \item \textbf{Foreningsmængde ($A \cup B$)} svarer til en bitvis \textbf{\texttt{OR}}-operation.
        \item \textbf{Fællesmængde ($A \cap B$)} svarer til en bitvis \textbf{\texttt{AND}}-operation.
        \item \textbf{Komplement ($A^c$)} svarer til en bitvis \textbf{\texttt{NOT}}-operation (enerkomplement).
        \item \textbf{Symmetrisk differens ($A \oplus B$)} svarer til en bitvis \textbf{\texttt{XOR}}-operation.
    \end{itemize}
\end{custombox}

% ---------- 04 Combinatorics & Probability ----------
\begin{definition}{Stokastisk eksperiment}
    Et stokastisk eksperiment er et eksperiment, der kan give forskellige resultater, selv om man gør alting på samme måde hver gang.
\end{definition}

\begin{definition}{Udfaldsrum, udfald og hændelser}
    \textbf{Udfaldsrum:} Mængden af alle mulige udfald af et stokastisk eksperiment kaldes udfaldsrummet og betegnes \( S \). Udfaldene kan være diskrete eller kontinuerte afhængigt af eksperimentets natur.
    
    \textbf{Udfald:} Et enkelt muligt resultat af et stokastisk eksperiment.
    
    \textbf{Hændelse:} En vilkårlig delmængde af udfaldsrummet, som kan bestå af et eller flere udfald.
\end{definition}

\begin{theorem}[Multiplikationsreglen]
    Lad en operation være beskrevet som en følge af $k$ trin. Antag, at følgende betingelser er opfyldt:
    
    \begin{itemize}
        \item Der er $n_1$ måder at udføre trin 1 på.
        \item Der er $n_2$ måder at udføre trin 2 på for hver måde at udføre trin 1 på.
        \item Der er $n_3$ måder at udføre trin 3 på for hver måde at udføre trin 2 på, og så videre.
    \end{itemize}
    
    Så er det samlede antal måder at udføre hele operationen på givet ved:
    
    \[
    n_1 \times n_2 \times \cdots \times n_k.
    \]
\end{theorem}

\begin{definition}[Optælling med og uden tilbagelægning]
    Når vi tæller antallet af måder at udvælge objekter fra en mængde på, er der to almindelige scenarier:

    \begin{itemize}
        \item \textbf{Med tilbagelægning}: Et objekt kan udvælges mere end én gang.
        \item \textbf{Uden tilbagelægning}: Når et objekt er udvalgt, kan det ikke vælges igen.
    \end{itemize}
\end{definition}

\begin{definition}[Permutation og kombination]
    \begin{itemize}
        \item \textbf{Permutation (rækkefølgen har betydning)}: Forskellige rækkefølger tælles som forskellige udfald, hvilket giver et større antal.
        \item \textbf{Kombination (rækkefølgen har ikke betydning)}: Rækkefølger behandles som ens, hvilket giver et mindre antal.
    \end{itemize}
\end{definition}

\begin{proposition}[Permutationer af \(n\) forskellige objekter]
    Antallet af ordnede arrangementer (permutationer) af $n$ forskellige objekter er
    \[
    n! = n \cdot (n-1) \cdot \dots \cdot 2 \cdot 1.
    \]
\end{proposition}

\begin{theorem}[Permutationer af delmængder]
    For heltal $n \ge r \ge 0$ er antallet af ordnede udvalg af $r$ forskellige objekter blandt $n$ forskellige objekter
    \[
    P_r^n = P(n,r) = n \cdot (n-1) \cdots (n-r+1) = \frac{n!}{(n-r)!}.
    \]
\end{theorem}

\begin{theorem}[Kombinationer]
    Antallet af kombinationer af $r$ elementer udvalgt fra en mængde med $n$ forskellige elementer er givet ved
    \[
    C_r^n = \binom{n}{r} = \frac{n!}{r!(n-r)!}
    \]
\end{theorem}

\begin{theorem}[Binomialsætningen]
    I algebra bruges binomialkoefficienten til at udvikle potenser af binomier. Ifølge binomialsætningen gælder

$$
(a+b)^n=\sum_{k=0}^n\binom{n}{k} a^k b^{n-k}
$$

\end{theorem}

\begin{proposition}[Med tilbagelægning]
For udvælgelser blandt $n$ typer gælder:
\begin{itemize}
    \item \textbf{Ordnet med tilbagelægning}: $n^r$ udfald.
    \item \textbf{Uordnet med tilbagelægning}: $\displaystyle \binom{n+r-1}{r}$ udfald. % (stars and bars)
\end{itemize}
\end{proposition}

\begin{axiom}[Sandsynlighedsregningens aksiomer]
    \begin{itemize}
        \item \textbf{Aksiom 1:} For enhver hændelse \(A\) gælder \( 0 \leq P(A) \leq 1 \).
        \item \textbf{Aksiom 2:} Sandsynligheden for udfaldsrummet \(S\) er \(P(S) = 1\).
        \item \textbf{Aksiom 3:} Hvis \(A_1, A_2, A_3, \cdots\) er disjunkte hændelser, så gælder \(P\left(A_1 \cup A_2 \cup A_3 \cdots\right) = P\left(A_1\right) + P\left(A_2\right) + P\left(A_3\right) + \cdots\)
    \end{itemize}
\end{axiom}

\begin{definition}[Sandsynligheden for en hændelse]
    Et sandsynlighedsmål tildeler et tal \(P(A)\in[0,1]\) til hver hændelse \(A\) i overensstemmelse med sandsynlighedsaksiomerne. I et endeligt udfaldsrum, hvor de elementære udfald er lige sandsynlige, reduceres dette til
    \[
    P(A) = \frac{\text{Antal gunstige udfald}}{\text{Samlet antal udfald}}
    \]
\end{definition}

\begin{theorem}{Regneregler for sandsynligheder}
    \begin{itemize}
        \item \textbf{Komplementreglen:} Sandsynligheden for komplementet til hændelse $A$ er
    
        \[
        P(A^c) = 1 - P(A)
        \]

        \item \textbf{Reglen for den tomme mængde:} Sandsynligheden for den tomme mængde er 0, dvs.
        
        \[P(\emptyset) = 0\]
        
        \item \textbf{Additionsreglen:} For to vilkårlige hændelser $A$ og $B$ er sandsynligheden for foreningsmængden af hændelserne $A$ og $B$ givet ved
        \[
        P(A \cup B) = P(A) + P(B) - P(A \cap B)
        \]

        \item \textbf{Differensreglen:} Sandsynligheden for differensen mellem hændelserne $A$ og $B$ er givet ved
        \[
        P(A-B)=P(A)-P(A \cap B)
        \]
        \item \textbf{Delmængdereglen:} Hvis $A$ er en delmængde af $B$ ($A \subset B$), så gælder
        \[P(A) \leq P(B).\]
    \end{itemize}
\end{theorem}

% ---------- 05 Conditional Probability ----------
\begin{definition}[Betinget sandsynlighed]
For to vilkårlige hændelser \( A \) og \( B \) med \( P(B) > 0 \) er den betingede sandsynlighed for \( A \) givet \( B \) defineret som:
\[
P(A \mid B) = \frac{P(A \cap B)}{P(B)}, \qquad P(B) > 0
\]
\end{definition}

\begin{axiom}[Aksiomer for betinget sandsynlighed]
    Fastlæg en hændelse $B$ med $P(B)>0$. Så er $A\mapsto P(A\mid B)$ et sandsynlighedsmål, og derfor gælder:
    \begin{itemize}
        \item \textbf{Aksiom 1:} For enhver hændelse \(A\) gælder \( 0 \leq P(A\mid B) \leq 1 \).
        \item \textbf{Aksiom 2:} $P(S\mid B)=1$ (og dermed $P(B\mid B)=1$).
        \item \textbf{Aksiom 3:} Hvis \(A_1, A_2, A_3, \cdots\) er disjunkte hændelser, så gælder
         
        \hspace*{2.8em}$P\left(A_1 \cup A_2 \cup A_3 \cdots \mid B\right)=P\left(A_1 \mid B\right)+P\left(A_2 \mid B\right)+P\left(A_3 \mid B\right)+\cdots$
    \end{itemize}
\end{axiom}

\begin{theorem}{Regler for betinget sandsynlighed}
    I de følgende regler antages $P(C)>0$.
    \begin{itemize}
        \item \textbf{Komplementreglen:} Sandsynligheden for komplementet til hændelse $A$ givet $C$ er
    
        \[
        P\left(A^c \mid C\right)=1-P(A \mid C)
        \]

        \item \textbf{Reglen for den tomme mængde:} Sandsynligheden for den tomme mængde givet en hændelse $C$ er 0, dvs.
        
        \[P(\emptyset \mid C) = 0\]
        
        \item \textbf{Additionsreglen:} For to vilkårlige hændelser $A$ og $B$ er sandsynligheden for foreningsmængden af hændelserne $A$ og $B$ givet en anden hændelse $C$
        \[
        P(A \cup B \mid C)=P(A \mid C)+P(B \mid C)-P(A \cap B \mid C)
        \]

        \item \textbf{Differensreglen:} Sandsynligheden for differensen mellem hændelserne $A$ og $B$ givet en anden hændelse $C$ er
        \[
            P(A-B \mid C)=P(A \mid C)-P(A \cap B \mid C)
        \]
        \item \textbf{Delmængdereglen:} Hvis $A \subset B$, så gælder
        \[
        P(A \mid C) \leq P(B \mid C)
        \]
    \end{itemize}
\end{theorem}

\begin{theorem}[Multiplikationsreglen]
    Hvis $P(A)>0$, så er $P(A\cap B)=P(A)P(B\mid A)$. Hvis $P(B)>0$, så er $P(A\cap B)=P(B)P(A\mid B)$. Når begge sandsynligheder er positive, gælder derfor:
    \[
    P(A \cap B) = P(A) P(B \mid A) = P(B) P(A \mid B)
    \]
\end{theorem}

\begin{theorem}[Loven om total sandsynlighed]
For en vilkårlig hændelse \(A\) og en hændelse \(B\) med $0<P(B)<1$ giver partitionen $B,B^c$:
    \begin{align*}
    P(A) &= P(A \cap B) + P(A \cap B^c) \\
    &= P(B) P(A \mid B) + P(B^c) P(A \mid B^c)
    \end{align*}

    Mere generelt gælder for enhver endelig partition \(B_1,\ldots,B_n\) af $S$ med $P(B_i)>0$ for alle $i$:
    \begin{align*}
    P(A) &= \sum_{i=1}^{n} P(A \mid B_i) P(B_i)
    \end{align*}
\end{theorem}

\begin{definition}[Uafhængighed]
    Hændelserne \(A\) og \(B\) er uafhængige, netop når
    \begin{align*}
        P(A \cap B) &= P(A)P(B).
    \end{align*}
    Hvis $P(B)>0$, er dette ensbetydende med $P(A\mid B)=P(A)$; hvis $P(A)>0$, er det ensbetydende med $P(B\mid A)=P(B)$.
\end{definition}

\begin{theorem}[Uafhængighed og De Morgans lov]
        Hvis \( A_1, A_2, \ldots, A_n \) er indbyrdes uafhængige, så gælder
        \begin{align*}
        &P\left(A_1 \cup A_2 \cup \cdots \cup A_n\right)=1-\left(1-P\left(A_1\right)\right)\left(1-P\left(A_2\right)\right) \cdots\left(1-P\left(A_n\right)\right)
        \end{align*}
\end{theorem}

\begin{definition}[Uafhængighed af flere hændelser]
    Hændelserne \(A_1,\ldots,A_k\) er \emph{indbyrdes uafhængige}, hvis der for enhver delmængde $J\subseteq\{1,\ldots,k\}$ af mindst to forskellige indeks gælder, at
    \[
      P\!\left(\bigcap_{j\in J}A_j\right)=\prod_{j\in J}P(A_j).
    \]
    Det er ikke tilstrækkeligt kun at kontrollere fællesmængden af alle $k$ hændelser.
\end{definition}

\begin{theorem}[Bayes' sætning]
    For hændelser $A$ og $B$ med $P(A)>0$ og $P(B)>0$ giver Bayes' sætning den første lighed nedenfor. Den udvidede nævner forudsætter desuden $0<P(B)<1$, så begge de viste betingede sandsynligheder er definerede:
    \begin{align*}
    P(B \mid A) = \frac{P(A \mid B) \cdot P(B)}{P(A)}
     = \frac{P(A \mid B) \cdot P(B)}{P(A \mid B) \cdot P(B) + P(A \mid B^c) \cdot P(B^c)}
    \end{align*}

    Hvis $B_1,B_2,\ldots$ udgør en endelig eller tællelig partition af $S$ med $P(B_i)>0$ for alle $i$, og $P(A)>0$, så gælder

    \begin{align*}
    P\left(B_j \mid A\right) = \frac{P\left(A \mid B_j\right) \cdot P\left(B_j\right)}{\sum_i P\left(A \mid B_i\right) \cdot P\left(B_i\right)}
    \end{align*}
\end{theorem}

% ---------- 06 Descriptive Statistics ----------

\begin{definition}[Population]
En \textbf{population} er den fuldstændige mængde af alle mulige observationer eller målinger, der er af interesse i en bestemt undersøgelse. I softwareudvikling kan det fx være alle mulige køretider for en algoritme, alle brugersessioner på et websted eller alle kodelinjer i et projekt.
\end{definition}

\begin{definition}[Stikprøve]
En \textbf{stikprøve} er en delmængde af populationen, som vi faktisk observerer eller måler. På grund af praktiske begrænsninger arbejder vi ofte med stikprøver frem for hele populationer.
\end{definition}

\begin{definition}[Kvalitative data]
\textbf{Kvalitative data} (også kaldet kategoriske data) består af ikke-numerisk information, der kan inddeles i kategorier. Eksempler er de programmeringssprog, der bruges i et projekt, brugernes vurdering af deres tilfredshed (tilfreds/neutral/utilfreds) eller fejlenes alvorlighedsgrad (kritisk/høj/middel/lav).
\end{definition}

\begin{definition}[Kvantitative data]
\textbf{Kvantitative data} består af numeriske målinger, der kan ordnes, og som man kan udføre matematiske operationer på. De kan yderligere opdeles i:
\begin{itemize}
    \item \textbf{Diskrete}: Tællelige værdier (antal fejl, kodelinjer, brugersessioner)
    \item \textbf{Kontinuerte}: Målbare værdier, der kan antage enhver værdi inden for et interval (svartider, hukommelsesforbrug, CPU-udnyttelse)
\end{itemize}
\end{definition}

\begin{definition}[Hyppighedsfordeling]
En \textbf{hyppighedsfordeling} er en tabel, der viser hyppigheden (antallet) af hver værdi eller klasse af værdier i et datasæt. Den kan angives som:
\begin{itemize}
    \item \textbf{Absolut hyppighed}: Det faktiske antal forekomster
    \item \textbf{Relativ hyppighed}: Andelen af det samlede antal observationer
    \item \textbf{Kumuleret hyppighed}: Den løbende sum af hyppighederne
\end{itemize}
\end{definition}

\begin{definition}[Aritmetisk gennemsnit]

For et datasæt med $n$ værdier $x_1, x_2, \ldots, x_n$ er det \textbf{aritmetiske gennemsnit} defineret som:

$$
\bar{x} = \frac{1}{n}\sum_{i=1}^{n} x_i = \frac{x_1 + x_2 + \cdots + x_n}{n}
$$

\end{definition}

\begin{definition}[Median]
    For et datasæt med $n$ værdier ordnet i stigende rækkefølge er medianen defineret som:
    \[
    \text{median} =
    \begin{cases}
        x_{\frac{n+1}{2}} & \text{hvis } n \text{ er ulige} \\
        \frac{x_{\frac{n}{2}} + x_{\frac{n}{2}+1}}{2} & \text{hvis } n \text{ er lige}
    \end{cases}
    \]
\end{definition}

\begin{definition}[Typetal]
\textbf{Typetallet} er den værdi, der forekommer hyppigst i et datasæt. Et datasæt kan have:
\begin{itemize}
    \item \textbf{Intet entydigt typetal}: Hvis alle observerede værdier har samme hyppighed, deler alle værdier den højeste hyppighed; efter nogle konventioner angives dette som »intet typetal«.
    \item \textbf{Ét typetal (unimodalt)}: Hvis én værdi forekommer hyppigst
    \item \textbf{Flere typetal (multimodalt)}: Hvis to eller flere værdier deler den højeste hyppighed
\end{itemize}
\end{definition}

\begin{definition}[Variationsbredde]
\textbf{Variationsbredden} for et datasæt er forskellen mellem den største og den mindste værdi:
\[
\text{Variationsbredde} = x_{\max} - x_{\min}
\]
\end{definition}

\begin{definition}[Stikprøvevarians og standardafvigelse]
For en stikprøve med $n$ værdier er \textbf{stikprøvevariansen}, betegnet $s^2$, defineret som:
$$
s^2 = \frac{1}{n-1}\sum_{i=1}^{n}(x_i - \bar{x})^2
$$
\textbf{Stikprøvestandardafvigelsen}, betegnet $s$, er kvadratroden af variansen:
$$
s = \sqrt{s^2} = \sqrt{\frac{1}{n-1}\sum_{i=1}^{n}(x_i - \bar{x})^2}
$$
\end{definition}

\begin{definition}[Variationskoefficient]
For data på forholdsskala med et stikprøvegennemsnit forskelligt fra nul er \textbf{variationskoefficienten} forholdet mellem stikprøvestandardafvigelsen og gennemsnittet:
\[
CV = \frac{s}{\bar{x}} \times 100\%
\]
\end{definition}

\begin{definition}[Percentil]
Den $p$'te \textbf{percentil} er en værdi, som cirka $p\%$ af de ordnede observationer ligger på eller under. Der findes flere interpolationskonventioner, så en numerisk percentilberegning bør angive, hvilken konvention der er brugt.
\end{definition}

\begin{definition}[Kvartiler]

Kvartiler deler et ordnet datasæt i fire lige store dele.

\begin{itemize}
    \item \textbf{Q1 (første kvartil)}: 25.-percentilen.
    \item \textbf{Q2 (anden kvartil)}: 50.-percentilen. Det er datasættets \textbf{median}.
    \item \textbf{Q3 (tredje kvartil)}: 75.-percentilen.
\end{itemize}
\end{definition}

\begin{definition}[Kvartilafstand]
    \textbf{Kvartilafstanden (IQR)} måler spredningen af de midterste 50\% af dataene og er defineret som
    $$
    IQR = Q_3 - Q_1,
    $$
    hvor $Q_1$ og $Q_3$ er første og tredje kvartil. En større IQR viser, at den centrale del af dataene er mere spredt, mens en mindre IQR viser, at den er tættere samlet.
\end{definition}

\begin{definition}[Normalfordeling]
En \textbf{normalfordeling} er en symmetrisk, klokkeformet fordeling, hvor middelværdien, medianen og typetallet alle er ens og ligger i centrum. Fordelingens form er fuldstændig bestemt af dens middelværdi ($\mu$) og standardafvigelse ($\sigma$).
\end{definition}

\begin{definition}[Skævhed]
    \textbf{Skævhed} måler asymmetrien i en datafordeling:
    \begin{itemize}
        \item \textbf{Positiv skævhed (højreskæv)}: Den længste hale strækker sig mod højre.
        \item \textbf{Negativ skævhed (venstreskæv)}: Den længste hale strækker sig mod venstre.
        \item \textbf{Symmetrisk}: De to sider har omtrent samme form.
    \end{itemize}
    Sammenligninger mellem gennemsnit og median kan være nyttige tommelfingerregler, men de definerer ikke skævhed.
\end{definition}

\begin{definition}[Boksplottets bestanddele]
Et standardboksplot efter Tukey viser kvartilerne og knurhårene:
\begin{itemize}
    \item boksen fra $Q_1$ til $Q_3$ og en linje ved medianen $Q_2$;
    \item knurhår, der strækker sig ud til de mest ekstreme observationer inden for $1.5\,IQR$ fra den pågældende kvartil; og
    \item observationer uden for knurhårene, som afsættes enkeltvis som mulige outliers.
\end{itemize}
\end{definition}

\begin{theorem}[Store tals lov]
    Hvis $X_1,X_2,\ldots$ er uafhængige og identisk fordelte stokastiske variable med endelig middelværdi $\mu$, så konvergerer deres stikprøvegennemsnit i sandsynlighed mod $\mu$, når stikprøvestørrelsen går mod uendelig.
\end{theorem}

\begin{theorem}[Den empiriske regel]
Hvis et datasæt er tilnærmelsesvis normalfordelt med stikprøvegennemsnit $\bar{x}$ og stikprøvestandardafvigelse $s$, så gælder:
\begin{itemize}
    \item Cirka \textbf{68.3\%} af observationerne ligger inden for 1 standardafvigelse fra gennemsnittet ($\bar{x} \pm s$).
    \item Cirka \textbf{95.4\%} af observationerne ligger inden for 2 standardafvigelser fra gennemsnittet ($\bar{x} \pm 2s$).
    \item Cirka \textbf{99.7\%} af observationerne ligger inden for 3 standardafvigelser fra gennemsnittet ($\bar{x} \pm 3s$).
\end{itemize}
\end{theorem}

% ---------- 07-09 Linear Algebra ----------

% 7

\begin{definition}{Lineære ligninger} En lineær ligning i variablene $x_1, x_2, \ldots, x_n$ er en ligning, der kan skrives på formen
\[
a_1 x_1 + a_2 x_2 + \cdots + a_n x_n = b
\]

hvor $b$ og koefficienterne $a_1, \ldots, a_n$ er konstanter.

\end{definition}

\begin{definition}{Lineære ligningssystemer} Et \textbf{lineært ligningssystem} (eller kort et \textbf{lineært system}) er en samling af en eller flere lineære ligninger med de samme variable. \end{definition}

\begin{definition}{Løsninger til et lineært ligningssystem} En \textbf{løsning} til et lineært ligningssystem i variablene $x_1, x_2, \ldots, x_n$ er en liste af tal $(s_1, s_2, \ldots, s_n)$, der opfylder alle systemets ligninger, når tallene indsættes på variablenes plads $x_1, x_2, \ldots, x_n$ i den nævnte rækkefølge. Mængden af alle mulige løsninger kaldes systemets \textbf{løsningsmængde}. To lineære ligningssystemer kaldes \textbf{ækvivalente}, hvis de har den samme løsningsmængde. \end{definition}

\begin{definition} En \textbf{matrix} er et rektangulært skema af tal ordnet i rækker og søjler. \textbf{Koefficientmatricen} for et lineært ligningssystem indeholder kun variablenes koefficienter, mens \textbf{totalmatricen} har en ekstra søjle med konstanterne fra ligningernes højre side. \end{definition}

\begin{definition} \textbf{Størrelsen} af en matrix er bestemt af antallet af rækker og søjler og angives som \emph{rækker} $\times$ \emph{søjler}. For eksempel har en koefficientmatrix med tre rækker og tre søjler størrelsen $3 \times 3$, og dens totalmatrix har størrelsen $3 \times 4$. \end{definition}

\begin{definition}{Echelonformer}
En matrix er på \textbf{echelonform}, hvis den opfylder følgende betingelser:
\begin{enumerate}
    \item Alle nulrækker står nederst.
    \item Hver ledende indgang i en række står i en søjle til højre for den ledende indgang i rækken ovenover.
    \item Alle indgange i en søjle under en ledende indgang er nul.
\end{enumerate}

Ud over echelonform kan en matrix også være på \textbf{reduceret echelonform} (RREF), som har følgende egenskaber:
\begin{enumerate}[resume]
    \item Den ledende indgang i hver række, der ikke er en nulrække, er 1.
    \item Hvert ledende 1-tal er den eneste indgang forskellig fra nul i sin søjle.
\end{enumerate}
\end{definition}

\begin{theorem}{Eksistenssætningen}
Et lineært ligningssystem er konsistent, hvis og kun hvis en echelonform af totalmatricen ikke har nogen række på formen
    \[
    \left[\begin{array}{llll}
    0 & \ldots & 0 & b
    \end{array}\right]
    \]
hvor $b$ er forskellig fra nul.
\end{theorem}

\begin{theorem}{Entydighed af den reducerede echelonform}
    Hver matrix er rækkeækvivalent med én og kun én matrix på reduceret echelonform
\end{theorem}

\begin{custombox}{Rækkereduktionsalgoritmen}
    Her beskriver vi en algoritme, der omdanner en vilkårlig matrix til en ækvivalent
    matrix på (reduceret) echelonform. Algoritmen er grundlaget for at løse lineære ligningssystemer ved hjælp af matricer.

    \begin{enumerate}
        \item Begynd med den søjle længst til venstre, der ikke er en nulsøjle. Det er en pivotsøjle, og pivotpositionen er øverst.
        \item Vælg en indgang forskellig fra nul i pivotsøjlen som pivot. Byt om nødvendigt rækker for at flytte denne indgang til pivotpositionen.
        \item Brug rækkeerstatninger til at skabe nuller i alle positioner under pivoten.
        \item Anvend trin 1-3 på delmatricen af alle indgange under og til højre for pivotpositionen. Gentag processen, indtil der ikke er flere rækker forskellige fra nul at ændre. (Nu har vi nået en echelonform af matricen.)
        \item Begynd med pivoten længst til højre, og arbejd opad og mod venstre: Skab nuller over hver pivot ved hjælp af rækkeoperationer. Hvis en pivot ikke er 1, gøres den til 1 ved en skalering. (Dette trin giver matricens reducerede echelonform.)
    \end{enumerate}

\end{custombox}

\begin{theorem}{Entydighedssætningen}
    Hvis et lineært ligningssystem er konsistent, indeholder løsningsmængden enten
    \begin{enumerate}[label=(\roman*)]
        \item en entydig løsning, når der ikke er nogen frie variable, eller
        \item uendeligt mange løsninger, når der er mindst én fri variabel.
    \end{enumerate}
\end{theorem}

% 8

\begin{definition}
    For ethvert positivt heltal $n$ lader vi $\mathbb{R}^n$ betegne samlingen af ordnede $n$-tupler, hvor hver indgang ligger i $\mathbb{R}$. Vi skriver ofte disse elementer som $n \times 1$-matricer. Vi definerer addition og multiplikation med skalar for vektorer i $\mathbb{R}^n$ på samme måde som for $\mathbb{R}^2$. Det vil sige, at vi regner koordinat for koordinat.
\end{definition}

\begin{definition}{Linearkombinationer}
    Givet vektorerne $\vect{v}_1, \vect{v}_2, \ldots, \vect{v}_p \in \mathbb{R}^n$ og skalarerne $c_1, c_2, \ldots, c_p \in \mathbb{R}$ kaldes vektoren $\vect{y}$ givet ved

\[
\vect{y} = c_1 \vect{v}_1 + \cdots + c_p \vect{v}_p
\]

en linearkombination af $\vect{v}_1, \vect{v}_2, \ldots, \vect{v}_p$ med vægtene $c_1, c_2, \ldots, c_p$.

\end{definition}

\begin{definition}{Udspændet af en mængde vektorer}
    Hvis $\vect{v}_1, \ldots, \vect{v}_p$ ligger i $\mathbb{R}^n$, betegnes mængden af alle linearkombinationer af $\vect{v}_1, \ldots, \vect{v}_p$ med \text{Span} $\left\{\vect{v}_1, \ldots, \vect{v}_p\right\}$, og den kaldes den delmængde af $\mathbb{R}^n$, der udspændes af $\vect{v}_1, \ldots, \vect{v}_p$. Med andre ord består udspændet af $\vect{v}_1, \ldots, \vect{v}_p$ af alle vektorer, der kan skrives på formen
    $$
    c_1 \vect{v}_1 + \cdots + c_p \vect{v}_p
    $$
    hvor $c_1, \ldots, c_p$ er skalarer.

\end{definition}

\begin{definition}{Matrixligning}
  Hvis $A$ er en $m \times n$-matrix med søjlerne $\vect{a}_1, \ldots, \vect{a}_n$, og hvis $\vect{x} \in \mathbb{R}^n$, så er produktet af $A$ og $\vect{x}$, betegnet $A \vect{x}$, givet ved

\[
A \vect{x} = \left[\begin{array}{llll}
\vect{a}_1 & \vect{a}_2 & \ldots & \vect{a}_n
\end{array}\right] \left[\begin{array}{c}
x_1 \\
x_2 \\
\vdots \\
x_n
\end{array}\right] = x_1 \vect{a}_1 + x_2 \vect{a}_2 + \cdots + x_n \vect{a}_n
\]
    
\end{definition}

\begin{definition}
    En indekseret mængde af vektorer $\left\{\vect{v}_1, \ldots, \vect{v}_p\right\} \subset \mathbb{R}^n$ siges at være \textbf{lineært uafhængig}, hvis vektorligningen
    
    \[
    x_1 \vect{v}_1 + \cdots + x_p \vect{v}_p = \vect{0}
    \]
    
    kun har den trivielle løsning, dvs. hvis den eneste løsning er $\left(x_1, \ldots, x_p\right) = (0, \ldots, 0)$. Tilsvarende siges mængden $\left\{\vect{v}_1, \ldots, \vect{v}_p\right\}$ at være \textbf{lineært afhængig}, hvis der findes vægte $c_1, \ldots, c_p$, som ikke alle er nul, så
    
    \[
    c_1 \vect{v}_1 + \cdots + c_p \vect{v}_p = \vect{0}
    \]
    
    Vi kalder en sådan ligning en \textit{lineær afhængighedsrelation}, når vægtene ikke alle er nul.
\end{definition}

\begin{theorem}
    Hvis $A$ er en $m \times n$-matrix med søjlerne $\vect{a}_1, \ldots, \vect{a}_n \in \mathbb{R}^m$, og hvis $\vect{b} \in \mathbb{R}^m$, så har matrixligningen

\[
A \vect{x} = \vect{b}
\]

samme løsningsmængde som vektorligningen

\[
x_1 \vect{a}_1 + x_2 \vect{a}_2 + \cdots + x_n \vect{a}_n = \vect{b}
\]

som igen har samme løsningsmængde som det lineære ligningssystem med totalmatricen

\[
\left[\begin{array}{lllll}
\vect{a}_1 & \vect{a}_2 & \ldots & \vect{a}_n & \vect{b}
\end{array}\right].
\]

\end{theorem}

\begin{theorem}
Lad \( A \) være en \( m \times n \)-matrix. Så er følgende udsagn ækvivalente:
\begin{enumerate}[label=(\alph*)]
    \item For hvert \( \vect{b} \in \mathbb{R}^m \) har ligningen \( A \vect{x} = \vect{b} \) en løsning.
    \item Hvert \( \vect{b} \in \mathbb{R}^m \) er en linearkombination af søjlerne i \( A \).
    \item Søjlerne i \( A \) udspænder \( \mathbb{R}^m \).
    \item Matricen \( A \) har en pivotposition i hver række.
\end{enumerate}
\end{theorem}

\begin{theorem}
Hvis \( A \) er en \( m \times n \)-matrix, \( \vect{u} \) og \( \vect{v} \) er vektorer i \( \mathbb{R}^n \), og \( c \) er en skalar, så gælder:

\begin{enumerate}[label=(\alph*)]
    \item \( A(\vect{u} + \vect{v}) = A \vect{u} + A \vect{v} \)  


    \item \( A(c \vect{u}) = c(A \vect{u}) \)  

\end{enumerate}
\end{theorem}

\begin{theorem}
    Antag, at $A \vect{x} = \vect{b}$ er konsistent for et bestemt $\vect{b}$, og lad $\vect{p}$ være en løsning. Så er løsningsmængden for $A \vect{x} = \vect{b}$ mængden af alle vektorer på formen
    \[
    \vect{w} = \vect{p} + \vect{v}_h,
    \]
    hvor $\vect{v}_h$ er en vilkårlig løsning til den homogene ligning $A \vect{x} = \vect{0}$.
\end{theorem}

\begin{theorem}
    \begin{enumerate}[(a)]
        \item Hvis en mængde vektorer indeholder nulvektoren, er mængden lineært afhængig.
        \item Hvis en mængde vektorer indeholder et skalarmultiplum af en anden vektor, er mængden lineært afhængig.
        \item Hvis en mængde vektorer indeholder flere vektorer, end der er indgange i hver vektor, er mængden lineært afhængig.
    \end{enumerate}
\end{theorem}

% 9

\begin{definition}{Matrixaddition}
    Lad $A$ og $B$ være $m \times n$-matricer. \textbf{Summen} af $A$ og $B$, betegnet $A + B$, er den $m \times n$-matrix, hvis indgange fås ved at lægge de tilsvarende indgange i $A$ og $B$ sammen.

    Givet matricerne

    $$
    A=\left[\begin{array}{cccc}
    a_{11} & a_{12} & \cdots & a_{1 n} \\
    a_{21} & a_{22} & \cdots & a_{2 n} \\
    \vdots & \vdots & \vdots & \vdots \\
    a_{m 1} & a_{m 2} & \cdots & a_{m n}
    \end{array}\right], \quad B=\left[\begin{array}{cccc}
    b_{11} & b_{12} & \cdots & b_{1 n} \\
    b_{21} & b_{22} & \cdots & b_{2 n} \\
    \vdots & \vdots & \vdots & \vdots \\
    b_{m 1} & b_{m 2} & \cdots & b_{m n}
    \end{array}\right]
    $$

    som begge har dimensionen $m \times n$, er summen $A+B$ altså defineret som

    $$
    A+B=\left[\begin{array}{cccc}
    a_{11}+b_{11} & a_{12}+b_{12} & \cdots & a_{1 n}+b_{1 n} \\
    a_{21}+b_{21} & a_{22}+b_{22} & \cdots & a_{2 n}+b_{2 n} \\
    \vdots & \vdots & \vdots & \vdots \\
    a_{m 1}+b_{m 1} & a_{m 2}+b_{m 2} & \cdots & a_{m n}+b_{m n}
    \end{array}\right]
    $$
\end{definition}

\begin{definition}{Multiplikation af en matrix med en skalar}
    Lad $A$ være en $m \times n$-matrix og $c$ en skalar. \textbf{Produktet} af $c$ og $A$, betegnet $cA$, er den $m \times n$-matrix, hvis indgange fås ved at gange hver indgang i $A$ med $c$. Det er altså defineret som

    $$
    c A=\left[\begin{array}{cccc}
    a_{11} & a_{12} & \cdots & a_{1 n} \\
    a_{21} & a_{22} & \cdots & a_{2 n} \\
    \vdots & \vdots & \vdots & \vdots \\
    a_{m 1} & a_{m 2} & \cdots & a_{m n}
    \end{array}\right] = \left[\begin{array}{cccc}
    c a_{11} & c a_{12} & \cdots & c a_{1 n} \\
    c a_{21} & c a_{22} & \cdots & c a_{2 n} \\
    \vdots & \vdots & \vdots & \vdots \\
    c a_{m 1} & c a_{m 2} & \cdots & c a_{m n}
    \end{array}\right]
    $$
    
\end{definition}

\begin{definition}
    For $A \in \mathbb{R}^{m \times n}$ og $B \in \mathbb{R}^{n \times p}$, hvor $B=\left[\begin{array}{lll}\vect{b_1} & \vect{b_2} \cdots & \vect{b_p}\end{array}\right]$, definerer vi produktet $A B$ ved formlen

\[
A B=\left[\begin{array}{llll}
A \vect{b_1} & A \vect{b_2} & \cdots & A \vect{b_p}
\end{array}\right]
\]

\end{definition}

\begin{definition}
    Lad $A$ være en kvadratisk matrix, dvs. $ A \in \mathbb{R}^{n \times n}$. Den $k$'te potens af $A$, betegnet $A^k$, er defineret som produktet af $A$ med sig selv $k$ gange. Det vil sige

    $$
    A^k=\underbrace{ AAA \cdot \cdots A}_{k \text { gange }}
    $$

    hvor $A$ optræder $k$ gange på højre side.
\end{definition}

\begin{definition}
    Givet en matrix $A \in \mathbb{R}^{m \times n}$ er den transponerede af $A$ den matrix $A^T$, hvis $i$'te søjle er den $i$'te række i $A$.
\end{definition}

\begin{definition}
    En kvadratisk matrix \( A \in \mathbb{R}^{n \times n} \) er invertibel (eller \textbf{ikke-singulær}), hvis der findes en matrix \( C \in \mathbb{R}^{n \times n} \), så \( A C = C A = I_n \). Matricen \( C \) kaldes den inverse til \( A \) og betegnes \( A^{-1} \). Altså er \( A^{-1} A = A A^{-1} = I_n \).
\end{definition}

\begin{theorem}

    Lad $A$, $B$, $C$ være matricer af samme størrelse, og lad $\alpha$, $\beta$ være skalarer. Så gælder
    \begin{enumerate}[(a)]
        \item $A+B=B+A$
        \item $(A+B)+C=A+(B+C)$
        \item $A+0=A$
        \item $\alpha(A+B)=\alpha A+\alpha B$
        \item $(\alpha+\beta) A=\alpha A+\beta A$
        \item $\alpha(\beta A)=(\alpha \beta) A$
    \end{enumerate}
\end{theorem}

\begin{theorem}
        Lad $A, B, C$ være matricer med passende dimensioner, og lad $\alpha$ være en skalar. Så gælder
\begin{enumerate}[(a)]
    \item $A(BC) = (AB)C$
    \item $A(B + C) = AB + AC$
    \item $(B + C)A = BA + CA$
    \item $\alpha(AB) = (\alpha A)B = A(\alpha B)$
    \item Hvis $A$ er $m\times n$, så er $I_mA=A=AI_n$.
\end{enumerate}
\end{theorem}

\begin{theorem}
    Lad $A$ og $B$ være matricer med passende dimensioner, og lad $\alpha$ være en skalar. Så gælder
    \begin{enumerate}[(a)]
        \item $(A^T)^T = A$
        \item $(A + B)^T = A^T + B^T$
        \item $(\alpha A)^T = \alpha A^T$
        \item $(AB)^T = B^T A^T$
    \end{enumerate}
\end{theorem}

\begin{theorem} Lad $A$ og $B$ være invertible $n \times n$-matricer. Så gælder:
\begin{enumerate}[(a)]
    \item $A^{-1}$ er invertibel, og
    \[
    \left(A^{-1}\right)^{-1}=A
    \]
    \item Produktet $A B$ er invertibelt, og
    \[
    (A B)^{-1}=B^{-1} A^{-1}
    \]
    \item Den transponerede af $A$ er også invertibel, dvs. $A^T$ er invertibel, og
    \[
    \left(A^T\right)^{-1}=\left(A^{-1}\right)^T
    \]
\end{enumerate}

\end{theorem}

\begin{theorem}
    Lad $A=\left[\begin{array}{ll}
    a & b \\
    c & d
    \end{array}\right]$. Hvis $ad-bc \neq 0$, er den inverse til $A$ givet ved

    \[
    A^{-1}=\frac{1}{ad-bc}\left[\begin{array}{cc}
    d & -b \\
    -c & a
    \end{array}\right]
    \]
\end{theorem}

\begin{theorem}
    En $n \times n$-matrix $A$ er invertibel, hvis og kun hvis $A$ er rækkeækvivalent med $I_n$. I så fald vil enhver følge af elementære rækkeoperationer, der reducerer $A$ til $I_n$, også omdanne $I_n$ til $A^{-1}$.
\end{theorem}

\begin{custombox}{Algoritme til at finde den inverse til en $n \times n$-matrix}
    Lad $A$ være en $n \times n$-matrix. Følg disse trin for at finde den inverse til $A$:
    \begin{enumerate}
        \item Opstil totalmatricen $[A | I_n]$.
        \item Udfør rækkeoperationer på $[A | I_n]$ for at reducere $A$ til $I_n$.
        \item Matricen i højre side af totalmatricen er $A^{-1}$.
    \end{enumerate}

    Hvis $A$ ikke er invertibel, vil algoritmen ikke kunne reducere $A$ til $I_n$.
\end{custombox}

\begin{theorem}
    Lad \( A \) være en invertibel matrix. Så har ligningen \( A \vect{x} = \vect{b} \) en entydig løsning, som er givet ved \( \vect{x} = A^{-1} \vect{b} \).
\end{theorem}

\begin{theorem}
    Lad \( A \) være en \( n \times n \)-matrix. Følgende udsagn er ækvivalente:
    \begin{enumerate}[(a)]
        \item $A$ er invertibel.
        \item $A$ er rækkeækvivalent med $I_n$.
        \item $A$ har $n$ pivotpositioner (dvs. én for hver række og søjle).
        \item Ligningen $A \vect{x}=\vect{0}$ har kun den trivielle løsning.
        \item Søjlerne i $A$ er lineært uafhængige.
        \item Ligningen \( A \vect{x} = \vect{b} \) har en entydig løsning for hvert \( \vect{b} \in \mathbb{R}^n \).
        \item Søjlerne i $A$ udspænder $\mathbb{R}^n$.
        \item \( \text{det}(A) \neq 0 \).
        \item Der findes en $n \times n$-matrix $C$, så $C A=I$.
        \item Der findes en $n \times n$-matrix $D$, så $A D=I$.
        \item $A^T$ er invertibel.
    \end{enumerate}
\end{theorem}

% ---------- 10-11 Calculus ----------

% 10

\begin{definition}[Grænseværdi]
Lad \( f(x) \) være en funktion, der er defineret i nærheden af et punkt \( c \). Vi siger, at \textbf{grænseværdien} af \( f(x) \) for \( x \) gående mod \( c \) er \( L \), hvilket skrives
\vspace{0.5em}
\[
\lim_{x \to c} f(x) = L
\]
\vspace{0.5em}
hvis vi kan få værdien af \( f(x) \) vilkårligt tæt på \( L \) ved at vælge et \( x \), der ligger tilstrækkelig tæt på \( c \), men ikke er lig med \( c \).
\end{definition}

\begin{definition}[Kontinuitet]
En funktion \( f \) er \textbf{kontinuert} i et punkt \( c \), hvis følgende tre betingelser er opfyldt:
\begin{enumerate}
    \item \( f(c) \) er defineret.
    \item \( \lim_{x \to c} f(x) \) eksisterer.
    \item \( \lim_{x \to c} f(x) = f(c) \).
\end{enumerate}
En funktion er kontinuert på et interval, hvis den er kontinuert i alle intervallets punkter.
\end{definition}

\begin{definition}[Den afledte funktion]
Den \textbf{afledte funktion} af en funktion \( f \) med hensyn til \( x \), betegnet \( f'(x) \), er funktionen
\[
f'(x) = \lim_{h \to 0} \frac{f(x+h) - f(x)}{h}
\]
forudsat at grænseværdien eksisterer. Hvis \( f'(x) \) eksisterer, siger vi, at \( f \) er \textbf{differentiabel} i \( x \).
\end{definition}

\begin{definition}
     For en funktion $y=f(x)$ kaldes de punkter på grafen, hvor grafen har hældningen nul, for stationære punkter. Med andre ord er de stationære punkter dér, hvor $f^{\prime}(x)=0$.
    \end{definition}

\begin{definition}[Den anden afledte]
Den \textbf{anden afledte} af en funktion \( f \), betegnet \( f''(x) \), er funktionen
\[
f''(x) = \frac{d}{dx} \left( f'(x) \right)
\]
forudsat at grænseværdien eksisterer.
\end{definition}

\begin{theorem}{Konstantreglen og potensreglen}
    \begin{enumerate}
        \item \textbf{Konstantreglen:} Hvis $c$ er en konstant, så er $\dfrac{d}{dx}(c) = 0$, \vspace{0.5em}
        \item \textbf{Potensreglen:} For ethvert reelt $n$ gælder i de punkter, hvor den reelle potensfunktion er defineret og differentiabel, at $\dfrac{d}{dx}x^n = n x^{n - 1}$. \vspace{0.5em}
    \end{enumerate}
\end{theorem}

\begin{theorem}{Skalarreglen og linearitetsreglen}
    \begin{enumerate}
        \setcounter{enumi}{2}
        \vspace{0.5em}
        \item \textbf{Skalarreglen:} \( \frac{d}{dx}[c \cdot f(x)] = c \cdot f'(x) \). \vspace{0.5em}
        \item \textbf{Linearitetsreglen:} \( \frac{d}{dx}[f(x) \pm g(x)] = f'(x) \pm g'(x) \). \vspace{0.5em}
    \end{enumerate}
\end{theorem}

\begin{theorem}{Produktreglen og kvotientreglen}
    \begin{enumerate}
        \setcounter{enumi}{4}
        \vspace{0.5em}
        \item \textbf{Produktreglen:} Hvis \(f(x)\) og \(g(x)\) er differentiable, så gælder \vspace{0.5em}
        \[
        \frac{d}{dx}\big(f(x)g(x)\big)=f'(x)g(x)+f(x)g'(x). \vspace{0.5em}
        \] \vspace{0.5em}
        \item \textbf{Kvotientreglen:} Hvis \(f(x)\) og \(g(x)\) er differentiable, og \(g(x)\neq0\), så gælder
        \vspace{0.5em}
        \[
        \frac{d}{dx}\left(\frac{f(x)}{g(x)}\right)=\frac{f'(x)g(x)-f(x)g'(x)}{[g(x)]^2}. \vspace{0.5em}
        \]
    \end{enumerate}
\end{theorem}

\begin{theorem}{Kædereglen}
    \vspace{0.5em}
    \begin{enumerate}
    \setcounter{enumi}{6}
        \item Hvis \( h(x) = f(g(x)) \) er en sammensat funktion, er dens afledte lig med den afledte af den ydre funktion (evalueret i den indre funktion) ganget med den afledte af den indre funktion. \vspace{0.5em}
        \[
        h'(x) = f'(g(x)) \cdot g'(x). \vspace{0.5em}
        \]
    \end{enumerate}
\end{theorem}

\begin{theorem}{Differentiationsregler}

    \begin{enumerate}
        \item \textbf{Konstantreglen:} Hvis $c$ er en konstant, så er $\dfrac{d}{dx}c=0$. \vspace{0.5em}
        \item \textbf{Potensreglen:} For ethvert reelt $n$ gælder i de punkter, hvor den reelle potensfunktion er defineret og differentiabel, at $\dfrac{d}{dx}x^n = n x^{n-1}$. \vspace{0.5em}
        \item \textbf{Skalarreglen:} Hvis $c$ er en konstant, så er $\dfrac{d}{dx}[c \cdot f(x)] = c \cdot f'(x)$. \vspace{0.5em}
        \item \textbf{Linearitet:} For differentiable $f,g$ og konstanter $a,b$ gælder $\dfrac{d}{dx}(a f + b g) = a f' + b g'$. \vspace{0.5em}
        \item \textbf{Produktreglen:} Hvis $f,g$ er differentiable, så er $\dfrac{d}{dx}(f g) = f'g + fg'$. \vspace{0.5em}
        \item \textbf{Kvotientreglen:} Hvis $f,g$ er differentiable, og $g(x)\neq0$, så er $\dfrac{d}{dx}\left(\frac{f}{g}\right) = \frac{f'g - fg'}{g^2}$. \vspace{0.5em}
        \item \textbf{Kædereglen:} Hvis $f,g$ er differentiable, så er $\dfrac{d}{dx}(f(g(x))) = f'(g(x))g'(x)$. \vspace{0.5em}
    \end{enumerate}

\end{theorem}

\begin{theorem}[Første afledtes test]
    Lad \(c\) være et kritisk punkt for en kontinuert funktion \(f\).
    \begin{itemize}
        \item Hvis \(f'(x)\) skifter fra negativ til positiv i \(c\), har \(f\) et \textbf{lokalt minimum} i \(c\).
        \item Hvis \(f'(x)\) skifter fra positiv til negativ i \(c\), har \(f\) et \textbf{lokalt maksimum} i \(c\).
        \item Hvis \(f'(x)\) ikke skifter fortegn i \(c\), har \(f\) intet lokalt ekstremum i \(c\). Det alene medfører ikke, at $c$ er et vendepunkt.
    \end{itemize}

\end{theorem}

\begin{theorem}[Anden afledtes test]
    Lad \(c\) være et stationært punkt for \(f\) (dvs. \(f'(c) = 0\)).
    \begin{itemize}
        \item Hvis \(f''(c) > 0\), har \(f\) et \textbf{lokalt minimum} i \(c\).
        \item Hvis \(f''(c) < 0\), har \(f\) et \textbf{lokalt maksimum} i \(c\).
        \item Hvis \(f''(c) = 0\), giver testen ingen konklusion.
    \end{itemize}
\end{theorem}

% 11

\begin{definition}[Funktion af to variable]
En \emph{funktion af to variable} er en regel, der til hvert ordnet par \((x, y)\) i en mængde \(D \subseteq \mathbb{R}^2\) knytter et entydigt reelt tal \(z\), som vi skriver som \(z = f(x, y)\). Mængden \(D\) kaldes \emph{definitionsmængden} for \(f\).

\emph{Værdimængden} for \(f\) er mængden af alle reelle tal \(z\), for hvilke der findes mindst ét \((x, y) \in D\), så \(f(x, y) = z\). Det er illustreret i \autoref{fig:ch11-domain-range}.
\end{definition}

\begin{definition}[Niveaukurve og konturkort]
En \textbf{niveaukurve} for en funktion \(f(x, y)\) er mængden af alle punkter \((x, y)\) i inputplanen, hvor funktionen har en konstant værdi, dvs. \(f(x, y) = c\) for en konstant \(c\).

En samling af niveaukurver for forskellige værdier af \(c\) udgør et \textbf{konturkort}.
\end{definition}

\begin{definition}[Partielle afledte]
Den \textbf{partielle afledte} af \( f(x, y) \) med hensyn til \( x \), betegnet \( \frac{\partial f}{\partial x} \) eller \( f_x \), er
$$
\frac{\partial f}{\partial x} = \lim_{h \to 0} \frac{f(x+h, y) - f(x, y)}{h}
$$
Den \textbf{partielle afledte} med hensyn til \( y \), betegnet \( \frac{\partial f}{\partial y} \) eller \( f_y \), er
$$
\frac{\partial f}{\partial y} = \lim_{h \to 0} \frac{f(x, y+h) - f(x, y)}{h}
$$
\end{definition}

\begin{definition}[Gradienten]
\textbf{Gradienten} af en funktion \( f(x, y) \), betegnet \( \nabla f \), er den vektorfunktion, der er defineret ved:

$$
\nabla f(x, y) = \left\langle \frac{\partial f}{\partial x}, \frac{\partial f}{\partial y} \right\rangle = \frac{\partial f}{\partial x} \mathbf{i} + \frac{\partial f}{\partial y} \mathbf{j}
$$

\end{definition}

% ---------- 12+ Number Theory and other universal concepts ----------

\begin{custombox}{{Kommutative, associative og distributive love}}


    \begin{itemize}
    \item Kommutativitet: $a+b=b+a$
    \item Associativitet: $(a+b)+c=a+(b+c)$
    \item Den distributive lov: $a \times(b+c)=a \times b+a \times c$
    \end{itemize}
    
\end{custombox}

\begin{custombox}{Lukkethed}
    \textbf{Addition:} Mængden af hele tal er lukket under addition. Det betyder, at hvis man tager to vilkårlige hele tal og lægger dem sammen, vil summen altid være et helt tal. Hvis \(a\) og \(b\) for eksempel er hele tal, så er \(a + b\) også et helt tal. Denne lukkethedsegenskab sikrer, at mængden af hele tal er stabil og fuldstændig under addition: Uanset hvor mange gange man lægger hele tal sammen, forbliver resultatet inden for mængden af hele tal.

\textbf{Multiplikation:} Mængden af hele tal er også lukket under multiplikation. Hvis man ganger to vilkårlige hele tal, vil produktet altid være et helt tal. Hvis \(a\) og \(b\) for eksempel er hele tal, så er \(a \times b\) også et helt tal. Denne egenskab garanterer, at multiplikation ligesom addition ikke giver resultater uden for mængden af hele tal, og bevarer dermed mængdens integritet under multiplikation.

\end{custombox}

\begin{custombox}{Ligheder og forskelle mellem rationale og irrationale tal}
    \begin{itemize}
    \item \textbf{Repræsentation:} Rationale tal kan skrives som brøker, mens irrationale tal ikke kan. Denne forskel gør irrationale tal mere komplicerede at håndtere, især i regneoperationer.

    \item \textbf{Decimaludvikling:} Decimaludviklingen af et rationalt tal er enten endelig eller periodisk, mens decimaludviklingen af et irrationalt tal er uendelig og ikke-periodisk.

    \item \textbf{Lukkethed:} De rationale tal er lukkede under addition, subtraktion, multiplikation og division (undtagen division med nul). De irrationale tal er ikke lukkede under disse operationer; for eksempel kan summen eller produktet af to irrationale tal undertiden være rationalt.

    \item \textbf{Tæthed:} De rationale tal ligger tæt på den reelle tallinje, dvs. at der mellem to vilkårlige rationale tal findes et andet rationalt tal. De irrationale tal ligger også tæt, men på en komplementær måde: De udfylder de »huller«, som de rationale tal efterlader, og sikrer dermed, at den reelle tallinje er sammenhængende uden afbrydelser.
\end{itemize}
\end{custombox}

\begin{theorem}{Aritmetikkens fundamentalsætning}
Ethvert helt tal større end 1 kan skrives entydigt som et produkt af primtal, bortset fra faktorernes rækkefølge.
\end{theorem}

% 13

\begin{custombox}{Primære booleske operationer}
    

\begin{itemize}
    \item \textbf{Komplementering} (betegnet \texttt{NOT}, fx $\overline{0}$)
    \item \textbf{Boolesk sum} (betegnet \texttt{OR}, fx $1 + 0$)
    \item \textbf{Boolesk produkt} (betegnet \texttt{AND}, fx $1 \cdot 0$)
\end{itemize}
\end{custombox}

\begin{custombox}{Grundlæggende booleske egenskaber}
\begin{itemize}
    \item \textbf{Dobbelt komplement:} \\
    \hspace*{15pt} $\overline{\overline{x}} = x$\hspace{15pt}
    
    \item \textbf{Idempotenslove:} \\
    \hspace*{15pt}$x + x = x$, \quad $x \cdot x = x$
    
    \item \textbf{Identitetslove:} \\
    \hspace*{15pt}$x + 0 = x$, \quad $x \cdot 1 = x$
    
    \item \textbf{Dominanslove:} \\
    \hspace*{15pt}$x + 1 = 1$, \quad $x \cdot 0 = 0$
    
    \item \textbf{Enhedsegenskaben:} \\
    \hspace*{15pt}$\overline{x} + x = 1$
    
    \item \textbf{Nulegenskaben:} \\
    \hspace*{15pt} $\overline{x} \cdot x = 0$
\end{itemize}
\end{custombox}

\begin{custombox}{Strukturelle booleske love}
\begin{itemize}
    \item \textbf{Kommutative love:} \\
    \hspace*{15pt}$x + y = y + x$, \quad $x \cdot y = y \cdot x$
    
    \item \textbf{Absorptionslove:} \\
    \hspace*{15pt}$x + (x \cdot y) = x$, \quad $x \cdot (x + y) = x$, \quad $\bar{x} \cdot \bar{y}+x=\bar{y}+x$
    
    \item \textbf{De Morgans love:} \\
    \hspace*{15pt} $\overline{x \cdot y} = \overline{x} + \overline{y}$, \quad $\overline{x + y} = \overline{x} \cdot \overline{y}$
    \item \textbf{Associative love:} \\
    \hspace*{15pt} $x + (y + z) = (x + y) + z$, \quad $x \cdot (y \cdot z) = (x \cdot y) \cdot z$
    
    \item \textbf{Distributive love:} \\
    \hspace*{15pt} $x \cdot (y + z) = (x \cdot y) + (x \cdot z)$, \quad $x + (y \cdot z) = (x + y) \cdot (x + z)$
\end{itemize}
\end{custombox}

\begin{custombox}{Avancerede booleske identiteter}
\begin{itemize}
    \item \textbf{Eksklusivt \texttt{OR} (\texttt{XOR}):} \\
    \hspace*{15pt} $x \oplus y = (x \cdot \overline{y}) + (\overline{x} \cdot y)$
    
    \item \textbf{Eksklusivt \texttt{NOR} (\texttt{XNOR}):} \\
    \hspace*{15pt} $\overline{x \oplus y} = (x \cdot y) + (\overline{x} \cdot \overline{y})$
    
    \item \textbf{Forsvindende modsætning:} \\
    \hspace*{15pt} $x + \overline{x} \cdot y = x + y$
    
    \item \textbf{FOIL-loven (First, Outer, Inner, Last -- første, ydre, indre, sidste):} \\
    \hspace*{15pt} $(x + y) \cdot (z + w) = (x \cdot z) + (x \cdot w) + (y \cdot z) + (y \cdot w)$
\end{itemize}
\end{custombox}

% 14

I de følgende definitioner er $f$ og $g$ ikke-negative funktioner, og $g(n)>0$ for alle tilstrækkelig store $n$.

\begin{definition}{Theta(\(\Theta\))}
    $f(n)$ er $\Theta(g(n))$, hvis der findes positive konstanter $c_1, c_2$ og $k$, så der for alle $n \geq k$ gælder
    \medskip
    \[
    c_1 \cdot g(n) \leq f(n) \leq c_2 \cdot g(n)
    \]
    \medskip
Det betyder, at $f(n)$ asymptotisk er begrænset både opadtil og nedadtil af $g(n)$, dvs. at $f(n)$ vokser med samme hastighed som $g(n)$, når $n$ bliver tilstrækkelig stor.
\end{definition}

\begin{proposition} 
    Lad $f(n)$ og $g(n)$ være to ikke-negative funktioner med $g(n)>0$ fra et vist trin, og antag, at der findes en konstant $c>0$, så
    \medskip
    \[
    \lim _{n \rightarrow \infty} \frac{f(n)}{g(n)}=c
    \]
    \medskip
    Så er $f(n)=\Theta(g(n))$. Hvis desuden
    \medskip
    \[
    \lim _{n \rightarrow \infty} \frac{f(n)}{g(n)}
    \]
    \medskip
    eksisterer, men \textbf{ikke} er lig med en positiv konstant, så er $f(n) \neq \Theta(g(n))$.
\end{proposition}

\begin{definition}{Store $\mathcal{O}$}
    $f(n)$ er $\mathcal{O}(g(n))$, hvis der findes positive konstanter $c$ og $k$, så der for alle $n \geq k$ gælder
    \medskip
    \[
    f(n) \leq c \cdot g(n)
    \]
    \medskip
    Det betyder, at $f(n)$ asymptotisk er begrænset opadtil af $g(n)$, dvs. at $f(n)$ højst vokser lige så hurtigt som $g(n)$, når $n$ bliver tilstrækkelig stor.
\end{definition}

\begin{definition}{Omega-notation}
    $f(n)$ er $\Omega(g(n))$, hvis der findes positive konstanter $c$ og $k$, så der for alle $n \geq k$ gælder
    \medskip
    \[
    f(n) \geq c \cdot g(n)
    \]
    \medskip
    Det betyder, at $f(n)$ asymptotisk er begrænset nedadtil af $g(n)$, dvs. at $f(n)$ mindst vokser lige så hurtigt som $g(n)$, når $n$ bliver tilstrækkelig stor.
\end{definition}

\begin{definition}{Lille $o$}
    $f(n)$ er $o(g(n))$, hvis der for enhver positiv konstant $\epsilon > 0$ findes en konstant $k \geq 0$, så der for alle $n \geq k$ gælder
    \medskip
    \[
    f(n) < \epsilon \cdot g(n)
    \]
    \medskip
    Det betyder, at $f(n)$ vokser strengt langsommere end $g(n)$, når $n$ bliver tilstrækkelig stor.
    
\end{definition}

\begin{proposition}
    Lad \( f(n) \) og \( g(n) \) være to ikke-negative funktioner med $g(n)>0$ fra et vist trin, og antag, at
    \medskip
    \[
    \lim_{n \rightarrow \infty} \frac{f(n)}{g(n)} = 0
    \]
    \medskip
    Så er \( f(n) = o(g(n)) \), dvs. at \( f(n) \) vokser strengt langsommere end \( g(n) \) for \( n \to \infty \).
\end{proposition}

\begin{definition}{Lille $\omega$}
    $f(n)$ er $\omega(g(n))$, hvis der for enhver positiv konstant $\epsilon > 0$ findes en konstant $k \geq 0$, så der for alle $n \geq k$ gælder
    \medskip
    \[
    f(n) > \epsilon \cdot g(n)
    \]
    \medskip
    Det betyder, at $f(n)$ vokser strengt hurtigere end $g(n)$, når $n$ bliver tilstrækkelig stor.
\end{definition}

\begin{proposition}
    Lad \( f(n) \) og \( g(n) \) være to ikke-negative funktioner med $g(n)>0$ fra et vist trin, og antag, at
    \medskip
    \[
    \lim_{n \rightarrow \infty} \frac{f(n)}{g(n)} = \infty
    \]
    \medskip
    Så er \( f(n) = \omega(g(n)) \), dvs. at \( f(n) \) vokser strengt hurtigere end \( g(n) \) for \( n \to \infty \).
\end{proposition}

        \begin{proposition}
            Lad \( f(n) \) og \( g(n) \) være to ikke-negative funktioner med $g(n)>0$ fra et vist trin. Hvis
            \medskip
            \[
            \limsup_{n \rightarrow \infty} \frac{f(n)}{g(n)} < \infty
            \]
            \medskip
            så er \( f(n) = \mathcal{O}(g(n)) \), dvs. at \( f(n) \) asymptotisk højst vokser lige så hurtigt som \( g(n) \).
        \end{proposition}

        \begin{proposition}
            Lad \( f(n) \) og \( g(n) \) være to ikke-negative funktioner med $g(n)>0$ fra et vist trin. Hvis
            \medskip
            \[
            \liminf_{n \rightarrow \infty} \frac{f(n)}{g(n)} > 0
            \]
            \medskip
            så er \( f(n) = \Omega(g(n)) \), dvs. at \( f(n) \) asymptotisk mindst vokser lige så hurtigt som \( g(n) \).
        \end{proposition}

        \begin{custombox}{Uformel grænseværdimetode}
            \[
                \begin{array}{ll}
                \displaystyle \lim_{n \to \infty} \frac{f(n)}{g(n)} = L<\infty & \Rightarrow f(n) = \mathcal{O}(g(n)) \\[10pt]
                \displaystyle \lim_{n \to \infty} \frac{f(n)}{g(n)} = L>0 & \Rightarrow f(n) = \Omega(g(n))
                \end{array}
                \]
                \medskip
                \textbf{Bemærk:} Denne metode er en uformel rettesnor og ikke en streng matematisk definition. Der findes tilfælde, hvor den grænseværdibaserede metode ikke kan anvendes, men den er ofte et bekvemt værktøj til at bestemme asymptotisk opførsel.
        \end{custombox}
